\examyear{2012}
\examera{平成24年}
\examfield{応用}
\examgroup{B}
\examnumber{16}
\examtitle{平成24年 B 第16問}
\examtag{遷移系}
\examtag{不動点}
\examtag{集合}
\examtag{論理}

遷移系 \((Q,R)\) とは，集合の組であって，\(R\subseteq Q\times Q\) を満たすものとする。
ただし，各 \(q\in Q\) に対し，\((q,q')\in R\) を満たす \(q'\) の個数は \(1\) 以上で有限であるとする。
また，\((q,q')\in R\) のかわりに \(q\to q'\) とも書く。
\(q=q_0\) から始まる無限遷移列 \(\pi:q_0\to q_1\to\cdots\) の集合を \(\Pi(q)\) と書く。
このとき，\(\pi(n)\) で \(q_n\) を表すことにする。\(Q\) の部分集合 \(X\) の関数
\[
  \AF(X)=\{q\in Q\mid \forall \pi\in\Pi(q)\ \exists n\ge0\ \pi(n)\in X\}
\]
を考える。

\begin{enumerate}
  \item 関数 \(F_X:2^Q\to 2^Q\) を
  \[
    F_X(Z)=X\cup \AX(Z)
  \]
  で定める。ただし
  \[
    \AX(Z)=\{q\in Q\mid \forall q'\ ((q,q')\in R\Rightarrow q'\in Z)\}
  \]
  とする。いま，集合の列 \(Z_n\) を，\(Z_0=X,\ Z_{n+1}=F_X(Z_n)\) で定義すれば，
  \[
    \AF(X)=\bigcup_{n=0}^{\infty}Z_n
  \]
  が成り立つことを示せ。
  \item \(\AF(X)\) は \(F_X\) の不動点であることを示せ。また，\(F_X\) の不動点の中で，包含関係について最小のものであることを示せ。
  ただし，\(F_X\) の不動点とは，\(F_X(Z)=Z\) を満たす集合 \(Z\) のことである。
\end{enumerate}